%! Compressing Images by the SVD — Unit 7, Lesson 7.2 — Slides (Beamer)
\documentclass[aspectratio=169, 11pt]{beamer}
\usepackage{linalg-beamer}   % provides \burgundyheader{} and \sectionlabel[color]{}
\newcommand{\uu}{\mathbf{u}}
\newcommand{\vv}{\mathbf{v}}
\newcommand{\T}{^{\top}}

\begin{document}

% TITLE SLIDE (CourseName is not defined in beamer, so the name is literal)
{
\setbeamercolor{background canvas}{bg=burgundy}
\begin{frame}[plain]
  \vspace{2.8cm}\hspace{0.55cm}%
  \begin{minipage}{0.9\textwidth}
    {\color{goldacc}\bfseries\small LESSON 7.2}\\[0.5cm]
    {\color{white}\bfseries\LARGE Compressing Images by the SVD}\\[1.0cm]
    {\color{blushmid}\small Linear Algebra~$\cdot$~Unit 7: The Singular Value Decomposition}
  \end{minipage}
\end{frame}
}

% HOOK
\begin{frame}[t, plain]
  \burgundyheader{A photo is a big matrix --- can we store less?}
  \sectionlabel{Hook}
  \begin{columns}[T]
    \begin{column}{0.52\textwidth}
      A $1000\times1000$ grayscale photo is a matrix of \textbf{one million} pixel values.\\[6pt]
      Lesson~7.1: the SVD splits any matrix into \emph{ordered layers}
      $\sigma_1\uu_1\vv_1\T,\ \sigma_2\uu_2\vv_2\T,\dots$
    \end{column}
    \begin{column}{0.44\textwidth}
      \begin{block}{The question}
        Keep only a few layers, throw the rest away --- \textbf{which} do you keep, and how much of the
        picture survives?
      \end{block}
    \end{column}
  \end{columns}
  \vspace{2pt}
  \begin{block}{Today}
    Build the layers, keep the biggest singular values, and count the storage saved.
  \end{block}
\end{frame}

% OUTER PRODUCT
\begin{frame}[t, plain]
  \burgundyheader{One new move: a column times a row}
  \sectionlabel[goldacc]{Idea}
  \begin{itemize}
    \setlength{\itemsep}{6pt}
    \item An \textbf{outer product} $\uu\vv\T$ is a column times a row --- and the result is a \emph{matrix}:
          \[ \begin{bmatrix} 1\\2 \end{bmatrix}\begin{bmatrix} 3 & 1 \end{bmatrix}
             =\begin{bmatrix} 3 & 1\\ 6 & 2 \end{bmatrix}. \]
    \item Every column is a multiple of $\uu$ --- the simplest nonzero matrix, called \textbf{rank-1}.
    \item It stores in only $m+n$ numbers (the entries of $\uu$ and $\vv$), not $m\times n$.
  \end{itemize}
\end{frame}

% SVD AS LAYERS
\begin{frame}[t, plain]
  \burgundyheader{The SVD is a sum of rank-1 layers}
  \sectionlabel{Skill}
  \begin{columns}[T]
    \begin{column}{0.5\textwidth}
      Multiply out $A=U\Sigma V\T$:
      \[ A=\sigma_1\uu_1\vv_1\T+\sigma_2\uu_2\vv_2\T+\cdots \]
      Each \textbf{layer} is a rank-1 outer product scaled by $\sigma_i$, ordered biggest $\sigma$ first.
    \end{column}
    \begin{column}{0.5\textwidth}
      \textbf{Spine} $A=\begin{bsmallmatrix} 5 & 3\\ 3 & 5 \end{bsmallmatrix}$, $\sigma_1=8,\sigma_2=2$
      (symmetric, so $\uu_i=\vv_i$):
      \[ 8\cdot\tfrac12\begin{bsmallmatrix} 1 & 1\\ 1 & 1 \end{bsmallmatrix}
         +2\cdot\tfrac12\begin{bsmallmatrix} 1 & -1\\ -1 & 1 \end{bsmallmatrix} \]
      \[ =\begin{bsmallmatrix} 4 & 4\\ 4 & 4 \end{bsmallmatrix}
         +\begin{bsmallmatrix} 1 & -1\\ -1 & 1 \end{bsmallmatrix}
         =\begin{bsmallmatrix} 5 & 3\\ 3 & 5 \end{bsmallmatrix}. \]
    \end{column}
  \end{columns}
\end{frame}

% COMPRESS
\begin{frame}[t, plain]
  \burgundyheader{Compress: keep the biggest layers}
  \sectionlabel{Skill}
  \begin{itemize}
    \setlength{\itemsep}{6pt}
    \item Keep the first $k$ layers: $A_k=\sigma_1\uu_1\vv_1\T+\cdots+\sigma_k\uu_k\vv_k\T$ --- the
          \textbf{closest} rank-$k$ matrix to $A$ (Eckart--Young).
    \item Spine: $A_1=\begin{bsmallmatrix} 4 & 4\\ 4 & 4 \end{bsmallmatrix}$, off by only $\pm1$.
    \item \textbf{Energy.} Total $=\sigma_1^2+\sigma_2^2=64+4=68$; layer~1 holds $64$, i.e.
          $64/68\approx\mathbf{94\%}$.
  \end{itemize}
  \vspace{2pt}
  \begin{block}{One layer already captures $94\%$ of the matrix}
    Big-$\sigma$ layers carry the most; the small tail is safe to drop.
  \end{block}
\end{frame}

% STORAGE + ROAD AHEAD
\begin{frame}[t, plain]
  \burgundyheader{The payoff --- and the road ahead}
  \sectionlabel[goldacc]{Payoff}
  \begin{itemize}
    \setlength{\itemsep}{6pt}
    \item \textbf{Storage.} A $1000\times1000$ image $=10^6$ numbers. Each layer costs $1+1000+1000=2001$.
          Keep $20$ layers: $\approx40{,}000$ numbers $\approx\mathbf{4\%}$ --- and it looks nearly identical.
    \item \textbf{Why the biggest $\sigma$'s?} They carry the most energy ($\sigma^2$), so dropping the tail
          loses the least. \emph{Every} matrix has this layered SVD, so every image compresses.
  \end{itemize}
  \vspace{2pt}
  \begin{block}{The road ahead}
    \textbf{7.3} principal component analysis (the same idea on \emph{data}) \;$\bullet$\; \textbf{7.4} the
    victory of orthogonality.
  \end{block}
\end{frame}

\end{document}
